\exercisesection{\S~9}

以下所考虑的所有 g-模都假设为有限维的。

\begin{enumerate}
\def\labelenumi{\arabic{enumi})}
\item
  若 \(m\) 是满足 \(\geq 0\) 的整数，则有 \(\dim \mathrm{E}(m\rho) = (m + 1)^{\mathrm{N}}\)，其中 \(\mathrm{N} = \operatorname{Card}(\mathrm{R}_{+})\)。
\item
  证明 \(d\) 上存在唯一的多项式函数 \(\mathfrak{h}^*\)，使得对所有 \(d(\lambda) = \dim \mathrm{E}(\lambda)\)，都有 \(\lambda \in \mathrm{P}_{++}\)；其次数为 \(\operatorname{Card}(\mathrm{R}_+)\)。有
\end{enumerate}

\[
d (w \lambda - \rho) = \varepsilon (w) d (\lambda - \rho) \quad \text { if } w \in W, \lambda \in \mathfrak {h} ^ {*}.
\]

特别地，函数 \(\lambda \mapsto d(\lambda - \rho)^2\) 在 W 下不变。推出 \(u\) 的中心中存在唯一元素 \(\mathrm{U}(\mathfrak{g})\)，使得对所有 \(\chi_{\lambda}(u) = d(\lambda)^2\)，都有 \(\lambda \in \mathfrak{h}^*\)（应用§8第5节的定理2）。当 \(\mathfrak{g} = \mathfrak{sl}(2,k)\) 时，有 \(u = C + 1\)，其中 \(C\) 是§1习题1中定义的元素。

\begin{enumerate}
\def\labelenumi{\arabic{enumi})}
\setcounter{enumi}{2}
\tightlist
\item
  令 \(k_{1}, \ldots, k_{l}\) 为 W 的不变量代数的特征次数（参见~第V章§5）。
\end{enumerate}

\begin{enumerate}
\def\labelenumi{\alph{enumi})}
\item
  证明对所有 \(j \geq 1\)，满足 \(k_{i} > j\) 的 i 的个数等于满足 \(\alpha \in R_{+}\) 的 \(\langle \rho, H_{\alpha} \rangle = j\) 的个数。（将情形化归为 R 不可约的情形，并使用第VI章§4习题6 c)。）（参见§5习题5 g)。
\item
  推出公式
\end{enumerate}

\[
\prod_ {\alpha \in \mathrm{R} _ {+}} \langle \rho , H _ {\alpha} \rangle = \prod_ {i = 1} ^ {l} (k _ {i} - 1)!.
\]

¶4) 假设 g 是单的，记 \(\gamma\) 为 \(R_{+}\) 中满足 \(H_{\gamma}\) 是 \(R^{\vee}\) 的最高根的元素；写作 \(H_{\gamma} = \sum_{\alpha \in B} n_{\alpha} H_{\alpha}\)。有 \(\langle \rho, H_{\gamma} \rangle = \sum n_{\alpha} = h - 1\)，其中 h 是 R 的 Coxeter 数（第VI章§1第11节命题31）。

\begin{enumerate}
\def\labelenumi{\alph{enumi})}
\tightlist
\item
  令 \(\alpha \in \mathrm{B}\)。证明对所有 \(\beta \in \mathrm{R}_{+}\)，
\end{enumerate}

\[
\langle \varpi_ {\alpha} + \rho , H _ {\beta} \rangle \leq h + n _ {\alpha} - 1,
\]

并且当 \(\beta = \gamma\) 时取等号。推出 \(\dim \mathrm{E}(\varpi_{\alpha})\) 的每个素因子都 \(\leq h + n_{\alpha} - 1\)。

\begin{enumerate}
\def\labelenumi{\alph{enumi})}
\setcounter{enumi}{1}
\tightlist
\item
  假设 \(\varpi_{\alpha}\) 不是极小权，即~\(n_{\alpha} \geq 2\)。令 \(m \in [2, n_{\alpha}]\) 且 \(p = h + m - 1\)。验证（参见~第VI章图版），存在 \(\beta \in \mathrm{R}_{+}\)，使得 \(\langle \varpi_{\alpha}, H_{\beta} \rangle = n_{\alpha}\)，且 \(\langle \rho, H_{\beta} \rangle = h - 1 - (n_{\alpha} - m)\)，从而 \(\langle \varpi_{\alpha} + \rho, H_{\beta} \rangle = p\)。推出，若 \(p\) 是素数，则 \(p\) 整除 \(\dim \mathrm{E}(\varpi_{\alpha})\)。（注意，对于 \(p\)，\(\langle \rho, H_{\beta} \rangle\) 不整除任何 \(\beta \in \mathrm{R}_{+}\)，参见~习题3。）
\end{enumerate}

\(c\) ) 当 \(\mathfrak{g}\) 为 \(\mathrm{G}_2\) 型（分别为 \(\mathrm{F}_4,\mathrm{E}_8\) 型）时，有 \(h = 6\)（分别为 12、30），且 \(\dim \mathrm{E}(\varpi_{\alpha})\) 被 7 整除（分别被 13、31 整除）；当 \(\mathfrak{g}\) 为 \(\mathrm{E}_6\) 型（分别为 \(\mathrm{E}_7\) 型），且 \(\varpi_{\alpha}\) 不是极小权时，\(\dim \mathrm{E}(\varpi_{\alpha})\) 被 13 整除（分别被 19 整除）。

¶ 5) a) 令 \(\alpha \in \mathrm{R}, x \in \mathfrak{g}^{\alpha}, y \in \mathfrak{g}^{-\alpha}\)，且 E 为 \(\mathfrak{g}\)-模。证明对所有 \(\lambda \in \mathrm{P}\)，

\[
\operatorname{Tr} ((x y) _ {\mathrm{E}} | \mathrm{E} ^ {\lambda}) = \operatorname{Tr} ((x y) _ {\mathrm{E}} | \mathrm{E} ^ {\lambda + \alpha}) + \lambda ([ x, y ]) \dim \mathrm{E} ^ {\lambda}.
\]

推出

\[
\sum_ {\lambda \in P} \lambda ([ x, y ]) \dim E ^ {\lambda}. e ^ {\lambda} = (1 - e ^ {- \alpha}) \sum_ {\lambda \in P} \operatorname{Tr} ((x y) _ {E} | E ^ {\lambda}). e ^ {\lambda}.
\]

\begin{enumerate}
\def\labelenumi{\alph{enumi})}
\setcounter{enumi}{1}
\tightlist
\item
  在 \(\mathfrak{h}^*\) 上取一个非退化的 W-不变对称双线性形式 \(\langle ,\rangle\)。令 \(\Delta\) 为向量空间 \(k[\mathrm{P}]\) 的自同态，使得对所有 \(\Delta(e^{\mu}) = \langle \mu ,\mu \rangle e^{\mu}\)，都有 \(\mu \in \mathrm{P}\)；若 \(a,b\in k[\mathrm{P}]\)，置
\end{enumerate}

\[
\Delta^ {\prime} (a, b) = \Delta (a b) - a \Delta (b) - b \Delta (a).
\]

证明

\[
\Delta (\mathrm{J} (e ^ {\mu})) = \langle \mu , \mu \rangle \mathrm{J} (e ^ {\mu}) \quad \text { for } \mu \in \mathrm{P},
\]

\[
\Delta^ {\prime} (e ^ {\lambda}, e ^ {\mu}) = 2 \langle \lambda , \mu \rangle e ^ {\lambda + \mu} \quad \mathrm{for} \lambda , \mu \in \mathrm{P},
\]

\[
\Delta^ {\prime} (a b, c) = a \Delta^ {\prime} (b, c) + b \Delta^ {\prime} (a, c) \quad \text { for } a, b, c \in k [ P ].
\]

\(c)\) 令 \(\lambda \in \mathrm{P}_{++}\)，\(c_{\lambda} = \operatorname{ch}(\operatorname{E}(\lambda))\) 且 \(d = \mathrm{J}(e^{\rho})\)。证明

\[
\Delta (c _ {\lambda} d) = \langle \lambda + \rho , \lambda + \rho \rangle c _ {\lambda} d.
\]

（使用 \(a\)、\(b\)、§6命题7的推论，以及第VI章§3第3节的公式 (3)。）

\(d)\) 由上述结果和§7第2节命题5 (iii) 推出 H. Weyl 公式的另一种证明。

\begin{enumerate}
\def\labelenumi{\alph{enumi})}
\setcounter{enumi}{4}
\tightlist
\item
  对所有 \(\lambda \in \mathfrak{h}^*\)，置 \(\dim \mathrm{E}^{\lambda} = m(\lambda)\)。由 \(a\) 推出
\end{enumerate}

\[
\begin{array}{c} \operatorname{Tr} ((x y) _ {\mathrm{E}} | \mathrm{E} ^ {\lambda}) = \sum_ {i = 0} ^ {+ \infty} (\lambda + i \alpha) ([ x, y ]) m (\lambda + i \alpha) \\ \sum_ {i = - \infty} ^ {+ \infty} (\lambda + i \alpha) ([ x, y ]) m (\lambda + i \alpha) = 0. \end{array}
\]

\(f)\) 令 \(\langle\cdot,\cdot\rangle\) 为 \(\mathfrak{g}\) 上的非退化不变对称双线性形式，其在 \(\mathfrak{h}\) 上的限制是上面所选形式的逆。令 \(\Gamma\) 为相应的 Casimir 元。假设 E 是单的；置 \(\Gamma_{\mathrm{E}}=\gamma.1\)，其中 \(\gamma\in k\)。利用 \(e\) 和§2第3节命题6，证明

\[
\gamma m (\lambda) = \langle \lambda , \lambda \rangle m (\lambda) + \sum_ {\alpha \in \mathrm{R}} \sum_ {i = 0} ^ {+ \infty} \langle \lambda + i \alpha , \alpha \rangle m (\lambda + i \alpha)
\]

对所有 \(\lambda \in \mathfrak{h}^*\)，并进而证明

\[
\begin{array}{l} \gamma m (\lambda) = \langle \lambda , \lambda \rangle m (\lambda) + \sum_ {\alpha \in \mathrm{R} _ {+}} m (\lambda) \langle \lambda , \alpha \rangle + 2 \sum_ {\alpha \in \mathrm{R} _ {+}} \sum_ {i = 1} ^ {+ \infty} m (\lambda + i \alpha) \langle \lambda + i \alpha , \alpha \rangle \\ = \langle \lambda , \lambda + 2 \rho \rangle m (\lambda) + 2 \sum_ {\alpha \in \mathrm{R} _ {+}} \sum_ {i = 1} ^ {+ \infty} m (\lambda + i \alpha) \langle \lambda + i \alpha , \alpha \rangle . \end{array}
\]

\(g)\) 继续假设 E 是单的；令 \(\omega\) 为其最高权。由 \(f)\) 推出，对所有 \(\lambda \in \mathfrak{h}^*\)，

\[
(\langle \omega + \rho , \omega + \rho \rangle - \langle \lambda + \rho , \lambda + \rho \rangle) m (\lambda) = 2 \sum_ {\alpha \in \mathrm{R} _ {+}} \sum_ {i = 1} ^ {+ \infty} m (\lambda + i \alpha) \langle \lambda + i \alpha , \alpha \rangle .
\]

（回忆由§7的命题5，若 \(\langle\omega+\rho,\omega+\rho\rangle>\langle\lambda+\rho,\lambda+\rho\rangle\) 是 E 的不同于 \(\lambda\) 的权，则 \(\omega\)。因此，上式给出了逐步计算 \(m(\lambda)\) 的方法。） \(^{18}\)

\begin{enumerate}
\def\labelenumi{\arabic{enumi})}
\setcounter{enumi}{5}
\tightlist
\item
  令 \(x \mapsto x^*\) 为 \(k[\mathrm{P}]\) 上的对合，将所有 \(e^p\) 的 \(e^{-p}\) 变为 \(p \in \mathrm{P}\)。\(a)\) 置 \(\mathrm{D} = d^* d = \prod_{\alpha \in \mathrm{R}} (1 - e^\alpha)\)。证明 \(d^* = (-1)^{\mathrm{N}} d\)，其中 \(\mathrm{N} = \mathrm{Card}(\mathrm{R}_+)\)，从而 \(\mathrm{D} = (-1)^{\mathrm{N}} d^2\)。
\end{enumerate}

\begin{enumerate}
\def\labelenumi{\alph{enumi})}
\setcounter{enumi}{1}
\tightlist
\item
  按下列公式在 \(\varepsilon\) 上定义两个线性形式 \(k[\mathrm{P}]\) 和 I：
\end{enumerate}

\[
\varepsilon (1) = 1, \quad \varepsilon (e ^ {p}) = 0 \quad \text { if } p \in \mathrm{P} - \{0 \}
\]

以及

\[
\mathrm{I} (f) = \frac {1}{m} \varepsilon (\mathrm{D}. f), \quad \mathrm{where} m = \mathrm{Card} (\mathrm{W}).
\]

利用公式 \(d = \mathrm{J}(e^{\rho})\)，证明 \(\mathrm{I}(1) = 1\)。

\begin{enumerate}
\def\labelenumi{\alph{enumi})}
\setcounter{enumi}{2}
\tightlist
\item
  令 \(\lambda \in \mathrm{P}_{++}\)，且 \(c_{\lambda} = \operatorname{ch} \mathrm{E}(\lambda) = \mathrm{J}(e^{\lambda + \rho}) / d\)。证明若 \(\mathrm{I}(c_{\lambda}) = 0\)，则 \(\lambda \neq 0\)。（方法同 \(b\)。）
\end{enumerate}

\(d\) ) 证明 I 在 \(\mathbf{Z}[\mathrm{P}]^{\mathrm{W}} = \operatorname {ch}\operatorname {R}(\mathfrak{g})\) 的子代数 \(k[\mathrm{P}]\) 上取整数值。若 E 是 \(\mathfrak{g}\)-模，则 E 中 \(\mathfrak{g}\) 的不变子空间的维数等于 I(ch E)。（将情形化归为 E 为单的情形，并使用 \(b\)）和 \(c\)。）\(e\) ) 有 \(\dim \mathrm{E} = \sum_{\lambda \in \mathrm{P}_{++}}\mathrm{I}(c_\lambda^*\mathrm{ch}\mathrm{E})d(\lambda)\)，其中 \(d(\lambda) = \dim \mathrm{E}(\lambda)\)。特别地：

\[
\mathrm{I} (c _ {\lambda} ^ {*} c _ {\mu}) = \delta_ {\lambda \mu} \quad \text { if } \lambda , \mu \in \mathrm{P} _ {+ +}.
\]

\(f)\) 若 \(\lambda, \mu, \nu \in \mathrm{P}_{++}\)，则命题2中的整数 \(m(\lambda, \mu, \nu)\) 等于 \(\mathrm{I}(c_{\lambda} c_{\mu} c_{\nu}^{*})\)。推出恒等式

\[
d (\lambda) d (\mu) = \sum_ {\nu \in \mathrm{P} _ {+ +}} m (\lambda , \mu , \nu) d (\nu) \quad \lambda , \mu , \nu \in \mathrm{P} _ {+ +}.
\]

（将 e) 应用于 g-模 \(\mathrm{E} = \mathrm{E}(\lambda) \otimes \mathrm{E}(\mu).\)

¶ 7) 沿用定理2证明中的记号。

\begin{enumerate}
\def\labelenumi{\alph{enumi})}
\tightlist
\item
  证明
\end{enumerate}

\[
f _ {\rho} (\mathrm{J} (e ^ {\mu})) = \prod_ {\alpha \in \mathrm{R} _ {+}} (e ^ {(\mu | \alpha) \mathrm{T} / 2} - e ^ {- (\mu | \alpha) \mathrm{T} / 2}).
\]

\begin{enumerate}
\def\labelenumi{\alph{enumi})}
\setcounter{enumi}{1}
\tightlist
\item
  取 \((\cdot |\cdot)\) 为典则双线性形式 \(\varPhi_{\mathrm{R}}\)（第VI章§1第12节）。证明
\end{enumerate}

\[
  f _ {\rho} (\mathrm{J} (e ^ {\mu})) \equiv d _ {\mu} \mathrm{T} ^ {\mathrm{N}} \left(1 + \frac {\mathrm{T} ^ {2}}{48} (\mu | \mu)\right) \pmod {\mathrm{T} ^ {\mathrm{N} + 3} \mathbf {R} [ [ \mathrm{T} ] ]},
\]

其中 \(d_{\mu} = \prod_{\alpha \in \mathrm{R}_{+}}(\mu |\alpha)\)。

\begin{enumerate}
\def\labelenumi{\alph{enumi})}
\setcounter{enumi}{2}
\tightlist
\item
  由 b) 以及等式 \(\mathrm{J}(e^{\lambda+\rho})=\mathrm{ch}(\mathrm{E}).\mathrm{J}(e^{\rho})\)，推出公式
\end{enumerate}

\[
  \sum_ {\mu \in \mathrm{P}} (\mu | \rho) ^ {2} \dim \mathrm{E} ^ {\mu} = \frac {\dim \mathrm{E}}{24} (\lambda | \lambda + 2 \rho).
\]

\(d\) ) 假设 \(\mathfrak{g}\) 是单的。证明 \((\rho |\rho) = \dim \mathfrak{g} / 24\)。（在 \(c\) 为 R 的最高根时应用 \(\lambda\)，并使用§6的习题3。）

¶8) 令 \(\psi\) 为 \(\mathfrak{h}^*\) 上次数为 \(r\) 的多项式函数。证明存在唯一的 \(\varPsi\) 上的多项式函数 \(\mathfrak{h}^*\)，它在 W 下不变，次数 \(\leq r\)，并且满足

\[
\sum_ {\mu \in \mathrm{P}} \psi (\mu) \dim \mathrm{E} ^ {\mu} = \varPsi (\lambda + \rho) \dim \mathrm{E}
\]

对任意最高权为 \(\lambda\) 的单 g-模 E。

（先处理 \(\psi(\mu)=(\mu|\nu)^{r}\) 的情形，其中 \(\nu\in P\) 不与任何根正交；为此使用定理2的证明中的同态 \(f_{\nu}\) 以及第V章§5第4节命题5 (i)。）

\begin{enumerate}
\def\labelenumi{\arabic{enumi})}
\setcounter{enumi}{8}
\tightlist
\item
  在 \(\mathfrak{g}\) 为 \(\mathrm{A}_2\) 型代数的情形，沿用第VI章图版I的记号。令 \(n, p\) 为满足 \(\geq 0\) 的整数。
\end{enumerate}

\[
a) \mathfrak {P} (n \alpha_ {1} + p \alpha_ {2}) = 1 + \inf (n, p).
\]

\begin{enumerate}
\def\labelenumi{\alph{enumi})}
\setcounter{enumi}{1}
\tightlist
\item
  令 \(\lambda = n\varpi_1 + p\varpi_2\)。则 \(\dim \mathrm{E}(\lambda) = \frac{1}{2}(n + 1)(p + 1)(n + p + 2)\)。若 \(\mathrm{E}(\lambda)\) 不是根权，即~\(\lambda\)，则 \(n \not\equiv p \pmod{3}\) 的权 0 的重数为 0；若 \(\lambda\) 是根权，则该重数为 \(1 + \inf(n,p)\)。
\end{enumerate}

\begin{enumerate}
\def\labelenumi{\arabic{enumi})}
\setcounter{enumi}{9}
\tightlist
\item
  在代数为 \(B_{2}\) 型的情形，沿用第VI章图版II的记号。令 n, p 为满足 \(\geq 0\) 的整数。
\end{enumerate}

\begin{enumerate}
\def\labelenumi{\alph{enumi})}
\tightlist
\item
  有
\end{enumerate}

\[
\begin{array}{c} \mathfrak {P} (n \alpha_ {1} + p (\alpha_ {1} + 2 \alpha_ {2})) = 1 + \frac {1}{2} p (p + 3) \\ \mathfrak {P} (n \alpha_ {2} + p (\alpha_ {1} + \alpha_ {2})) = [ p ^ {2} / 4 ] + p + 1 \\ \mathfrak {P} (n (\alpha_ {1} + \alpha_ {2}) + p (\alpha_ {1} + 2 \alpha_ {2})) = [ n ^ {2} / 4 ] + n + 1 + n p + \frac {1}{2} p (p + 3). \end{array}
\]

\begin{enumerate}
\def\labelenumi{\alph{enumi})}
\setcounter{enumi}{1}
\tightlist
\item
  令 \(\lambda = n(\alpha_{1} + \alpha_{2}) + p(\alpha_{1} + 2\alpha_{2})\)。\(\operatorname{E}(\lambda)\) 中权 0 的重数为
\end{enumerate}

\[
[ n / 2 ] + 1 + n p + p.
\]

\begin{enumerate}
\def\labelenumi{\arabic{enumi})}
\setcounter{enumi}{10}
\item
  假设 \(\mathfrak{g}\) 不是秩为 1 的代数的乘积。令 \(n\in \mathbf{N}\)。证明存在一个单 \(\mathfrak{g}\)-模，其某个权的重数 \(\geq n\)。（反设不然。令 \(E_{\lambda}\) 为最高权 \(\lambda \in P_{++}\) 的单 g-模。当 \(\dim E_{\lambda}\)\(E_{\lambda}\) 时，比较 \((\lambda|\lambda) \to \infty\) 的维数和其不同权的数目（使用定理2的记号）。）
\item
  令 \(U^{0}\) 为 \(\mathrm{U}(\mathfrak{g})\) 中 h 的交换子。若 g 的秩为 1，则 \(U^{0}\) 是交换的。若 g 的秩 \(\geq 2\)，则 \(U^{0}\) 有任意大的有限维单表示。（使用习题11和§7习题23 a)。）
\item
  令 R 为向量空间 V 中的根系。若 R 是两个根系 \(v_{1}, v_{2}\) 和 \(R_{1}\) 的直和（第VI章§1第2节），且 \(R_{2}\) 属于 Rᵢ 生成的 V 的向量子空间（\(v_{i}\)），则称 V 的两个元素 \(R_{i}, i = 1, 2\) 不相交。证明，属于 R 的同一室的 \(V_{R}\) 中两个元素不相交，当且仅当它们正交。
\end{enumerate}

¶ 14) a) 令 \(\mu, \nu \in \mathrm{P}_{++}\)，令 \(\gamma\) 为 \(\mathrm{E}(\mu)\) 的一个权。令 \(\rho_{\mu}\) 为 \(\mathfrak{g}\) 在 \(\mathrm{E}(\mu)\) 上的表示。令 \(X_{\alpha} \in \mathfrak{g}^{\alpha} - \{0\}, Y_{\alpha} \in \mathfrak{g}^{-\alpha} - \{0\}\)。若 \(\alpha \in \mathrm{B}\)，置 \(\nu_{\alpha} = \nu(H_{\alpha})\)，并

\[
\mathrm{E} ^ {+} (\mu , \gamma , \nu) = \mathrm{E} (\mu) ^ {\gamma} \cap \bigcap_ {\alpha \in \mathrm{B}} \mathrm{Ker} \rho_ {\mu} (X _ {\alpha}) ^ {\nu_ {\alpha} + 1},
\]

\[
\mathrm{E} ^ {-} (\mu , \gamma , \nu) = \mathrm{E} (\mu) ^ {\gamma} \cap \bigcap_ {\alpha \in \mathrm{B}} \operatorname{Ker} \rho_ {\mu} (Y _ {\alpha}) ^ {\nu_ {\alpha} + 1},
\]
% semantic-fragments: legacy-input-seam
\relax{}
\[
d ^ {+} (\mu , \gamma , \nu) = \dim \mathrm{E} ^ {+} (\mu , \gamma , \nu), \quad d ^ {-} (\mu , \gamma , \nu) = \dim \mathrm{E} ^ {-} (\mu , \gamma , \nu).
\]

对所有 \(\lambda \in \mathfrak{h}^*\)，置 \(\lambda^{*} = -w_{0}\lambda\)，其中 \(w_{0}\) 是 W 中把 B 变为 \(-B\) 的元素。证明

\[
d ^ {+} (\mu , \gamma , \nu) = d ^ {-} (\mu , - \gamma^ {*}, \nu^ {*}).
\]

\begin{enumerate}
\def\labelenumi{\alph{enumi})}
\setcounter{enumi}{1}
\item
  令 \(\lambda_1, \lambda_2 \in \mathrm{P}_{++}\)，V 为 \(\mathfrak{g}\)-模 \(\mathrm{Hom}_k(\mathrm{E}(\lambda_1^*), \mathrm{E}(\lambda_2))\)，U 为由 \(\varphi \in V\) 中满足 \(Y_\alpha. \varphi = 0\)（对所有 \(\alpha \in B\)）者组成的集合，\(w\) 为 \(\mathrm{E}(\lambda_1^*)\) 中的本原向量。证明 \(\varphi \mapsto \varphi(w)\) 是从 U 到由 \(v \in \mathrm{E}(\lambda_2)\) 中满足 \(Y_\alpha^{\lambda_1^*(H_\alpha)+1}. v = 0\)（对所有 \(\alpha \in B\)）者组成的集合的同构。（为证满射，使用§7的习题15。）
\item
  采用命题2的记号，证明对 \(\lambda_1, \lambda_2, \lambda \in \mathrm{P}_{++}\)，
\end{enumerate}

\[
\begin{array}{r} m (\lambda_ {1}, \lambda_ {2}, \lambda) = d ^ {+} (\lambda , \lambda_ {2} - \lambda_ {1} ^ {*}, \lambda_ {1} ^ {*}) = d ^ {-} (\lambda , \lambda_ {1} - \lambda_ {2} ^ {*}, \lambda_ {1}) \\ = d ^ {+} (\lambda_ {1}, \lambda - \lambda_ {2}, \lambda_ {2}) = d ^ {-} (\lambda_ {1}, \lambda_ {2} ^ {*} - \lambda^ {*}, \lambda_ {2} ^ {*}). \end{array}
\]

（注意 \(m(\lambda_1,\lambda_2,\lambda)\) 是 \(\mathfrak{g}\)-不变元在

\[
\mathrm{E} (\lambda) ^ {*} \otimes \mathrm{E} (\lambda_ {1}) \otimes \mathrm{E} (\lambda_ {2}),
\]

中所成空间的维数，因而等于 \(m(\lambda_1^*, \lambda, \lambda_2)\)。）

\(d)\) 令 \(\lambda_1, \lambda_2 \in \mathrm{P}_{++}\)，并令 \(\lambda\) 为 \(\mathrm{P}_{++} \cap \mathrm{W}.\) ( \(\lambda_1 - \lambda_2^*\) ) 的唯一元素。我们有

\[
m (\lambda_ {1}, \lambda_ {2}, \lambda) = 1.
\]

（使用 c)。）由此推出 \(\mathrm{E}(\lambda_1) \otimes \mathrm{E}(\lambda_2)\) 含有唯一同构于 \(\mathrm{E}(\lambda)\) 的子模。

\begin{enumerate}
\def\labelenumi{\alph{enumi})}
\setcounter{enumi}{4}
\tightlist
\item
  保留 \(d\) ) 的记号。证明下列条件等价：
\end{enumerate}

\begin{enumerate}
\def\labelenumi{(\roman{enumi})}
\item
  \(\lambda = \lambda_{1} + \lambda_{2}\)
\item
  \(\| \lambda \| = \| \lambda_1 + \lambda_2 \|\)
\item
  \(\lambda_{1}\) 与 \(\lambda_2^*\) 正交
\item
  \(\lambda_{1}\) 与 \(\lambda_2^*\) 不相交（习题13）
\item
  \(\lambda_{1}\) 与 \(\lambda_{2}\) 不相交。
\end{enumerate}

由此断定，\(\mathrm{E}(\lambda_{1})\otimes\mathrm{E}(\lambda_{2})\) 不是单模，除非 \(\lambda_{1}\) 与 \(\lambda_{2}\) 不相交（从而给出§7习题18 f) 的又一个证明）。

\begin{enumerate}
\def\labelenumi{\arabic{enumi})}
\setcounter{enumi}{14}
\tightlist
\item
  置 \(\mathrm{N} = \mathrm{Card}(\mathrm{R}_{+})\)，\(c = \prod_{\alpha \in \mathrm{R}_{+}} \langle \rho, H_{\alpha} \rangle\)，并置 \(d(\lambda) = \dim \mathrm{E}(\lambda)\)（对 \(\lambda \in \mathrm{P}_{++}\)）。证明对任意实数 \(s > 0\)，
\end{enumerate}

\[
\sum_ {\lambda \in \mathrm{P} _ {+ +}} d (\lambda) ^ {- s} \leq \frac {1}{c} \left(\sum_ {m = 1} ^ {\infty} m ^ {- s}\right) ^ {\mathrm{N}}.
\]

由此推出 \(\sum_{\lambda \in \mathrm{P}_{+ + }}d(\lambda)^{-s} <   + \infty\)，若 \(s > 1\)

¶ 16) a) 取 g 为 \(F_{4}\) 型，并采用第VI章图版VIII的记号。当 i = 1, 2, 3, 4 时，置

\[
\mathcal {X} _ {i} = \mathrm{P} _ {+ +} \cap (\varpi_ {i} - \mathrm{Q} _ {+}).
\]

\(\mathrm{E}(\varpi_i)\) 的权集合是诸 \(\mathrm{W}\omega\) 的不交并，其中 \(\omega\) 属于 \(\mathcal{X}_i\)（参见§7命题5 (iv)）。我们有：

\[
\mathcal {X} _ {1} = \{0, \varpi_ {1}, \varpi_ {4} \};
\]

\[
\mathcal {X} _ {2} = \{0, \varpi_ {1}, \varpi_ {2}, \varpi_ {3}, \varpi_ {4}, \varpi_ {1} + \varpi_ {4}, 2 \varpi_ {4} \};
\]

\[
\mathcal {X} _ {3} = \{0, \varpi_ {1}, \varpi_ {3}, \varpi_ {4} \};
\]

\[
\mathcal {X} _ {4} = \{0, \varpi_ {4} \}.
\]

\begin{enumerate}
\def\labelenumi{\alph{enumi})}
\setcounter{enumi}{1}
\tightlist
\item
  借助定理2证明：
\end{enumerate}

\(\dim \mathrm{E}(\varpi_1) = 52,\quad \dim \mathrm{E}(\varpi_2) = 1274,\quad \dim \mathrm{E}(\varpi_3) = 273,\quad \dim \mathrm{E}(\varpi_4) = 26.\)

\begin{enumerate}
\def\labelenumi{\alph{enumi})}
\setcounter{enumi}{2}
\tightlist
\item
  利用第V章§3的命题1以及第VI章的图版，证明
\end{enumerate}

\[
\mathrm{Card} (\mathrm{W} \varpi_ {1}) = 2 ^ {7} 3 ^ {2} 2 ^ {- 3} (3!) ^ {- 1} = 24.
\]

类似地计算 \(\mathrm{Card}(\mathrm{W}\varpi_{2}),\ldots,\mathrm{Card}(\mathrm{W}.2\varpi_{4})\)。由此推出 \(\mathrm{E}(\varpi_{2})\) 的权数为 553；由于该数严格小于 \(\dim\mathrm{E}(\varpi_{2})\)，这些权中有一个的重数 \(\geq2\)。

\begin{enumerate}
\def\labelenumi{\alph{enumi})}
\setcounter{enumi}{3}
\item
  对 \(\varpi_{1},\varpi_{3},\varpi_{4}\) 作类似的计算。利用§7的习题21推出：若 \(\rho\) 是 g 的非零单表示，则 \(\rho\) 具有一个重数 \(\geq2\) 的权。
\item
  对 \(\mathrm{E}_8\) 型的单代数证明同样的结果。
\end{enumerate}

\(f)\) 令 \(\mathfrak{a}\) 为可分裂单李代数。由 \(d\) )、\(e)\) 以及§7的命题7和命题8，推出下列性质等价：

\begin{enumerate}
\def\labelenumi{(\roman{enumi})}
\item
  \(\mathfrak{a}\) 具有一个非零单表示，其所有权的重数皆为 1；
\item
  \(\mathfrak{a}\) 既非 \(\mathrm{F}_4\) 型也非 \(\mathrm{E}_8\) 型。
\end{enumerate}

